\begin{figure}[htbp]
\centering
\begin{tikzpicture}[
  >=Stealth,
  every node/.style={font=\small},
  first/.style={very thick,blue!65!black},
  second/.style={very thick,orange!70!black},
  port/.style={circle,fill=gray!70,inner sep=1.4pt},
  retained/.style={circle,draw,fill=white,inner sep=1.7pt}
]
\begin{scope}
\node at (0.9,2.05) {idempotent $e\in\D{4}$};
\foreach \i/\y in {1/1.2,2/0.4,3/-0.4,4/-1.2} {
  \coordinate (eL\i) at (0,\y);
  \coordinate (eR\i) at (1.8,\y);
  \node[anchor=south] at ([yshift=2pt]eL\i) {$\i$};
  \node[anchor=south] at ([yshift=2pt]eR\i) {$\i$};
}
\draw[first] (eL1)--(eR1);
\draw[second] (eL2)--(eR4);
\draw[thick,gray] (eL3) .. controls (-0.55,-0.4) and (-0.55,-1.2) .. (eL4);
\draw[thick,gray] (eR2) .. controls (2.35,0.4) and (2.35,-0.4) .. (eR3);
\foreach \p in {eL3,eL4,eR2,eR3} \node[port] at (\p) {};
\foreach \p in {eL1,eL2,eR1,eR4} \node[retained] at (\p) {};
\node[align=center] at (0.9,-1.85)
  {$X=\{1,2\}$,\quad $Y=\{1,4\}$\\$p(1)=1,\quad p(2)=4$};
\end{scope}

\draw[->,thick] (2.7,0)--node[above] {glue} (3.4,0);

\begin{scope}[shift={(4.3,0)}]
\node at (0.9,2.05) {connector in $e\id_4e$};
\foreach \i/\y in {1/1.2,2/0.4,3/-0.4,4/-1.2} {
  \coordinate (cR\i) at (0,\y);
  \coordinate (cL\i) at (1.8,\y);
  \node[anchor=south] at ([yshift=2pt]cR\i) {$R_{\i}$};
  \node[anchor=south] at ([yshift=2pt]cL\i) {$L_{\i}$};
}
\draw[first] (cR1)--(cL1);
\draw[second] (cR4)--(cL4);
\draw[second] (cL4) .. controls (2.35,-1.2) and (2.35,-0.4) .. (cL3);
\draw[second] (cL3)--(cR3);
\draw[second] (cR3) .. controls (-0.55,-0.4) and (-0.55,0.4) .. (cR2);
\draw[second] (cR2)--(cL2);
\foreach \p in {cR2,cR3,cL3,cL4} \node[port] at (\p) {};
\foreach \p in {cR1,cR4,cL1,cL2} \node[retained] at (\p) {};
\node at (0.9,-1.85) {$\gamma(1)=1,\quad\gamma(4)=2$};
\end{scope}

\draw[->,thick] (7.0,0)--node[above,align=center] {retain and\\renumber} (8.15,0);

\begin{scope}[shift={(9.1,0)}]
\node at (0.9,2.05) {$R_e(e)=\id_2$};
\coordinate (dL1) at (0,0.6);
\coordinate (dL2) at (0,-0.6);
\coordinate (dR1) at (1.8,0.6);
\coordinate (dR2) at (1.8,-0.6);
\draw[first] (dL1)--(dR1);
\draw[second] (dL2)--(dR2);
\foreach \p in {dL1,dL2,dR1,dR2} \node[retained] at (\p) {};
\node[anchor=south] at ([yshift=2pt]dL1) {$1$};
\node[anchor=south] at ([yshift=2pt]dL2) {$2$};
\node[anchor=south] at ([yshift=2pt]dR1) {$1$};
\node[anchor=south] at ([yshift=2pt]dR2) {$2$};
\node[align=center] at (0.9,-1.85)
  {$x$ and $p(x)$ share a number;\\right $4=p(2)$ becomes $2$};
\end{scope}
\end{tikzpicture}
\caption{A rank-$2$ idempotent of degree $4$ and its reduction. In the
left panel, the gray caps are prescribed and the colored transversals
are retained. The connector paths are $R_1$--$L_1$ and
$R_4$--$L_4$--$L_3$--$R_3$--$R_2$--$L_2$, giving
$\gamma=p^{-1}$. Numbering left $x$ and right $p(x)$ alike yields
$R_e(e)=\id_2$. For a general corner element $z$, only the prescribed
outer caps are deleted; every pair among retained ports, including
caps, remains.}
\label{fig:corner-reduction}
\end{figure}
